\section{A fourth-order rotation estimate}
\label{sec:rotation}

The analytic input is a finite-graph inequality. Rotations about two
different axes produce a second-order effective rotation about a third
axis. A fourth variation of the partition function detects this effect;
after summing a family of test functions, its potentially large terms
form a nonnegative square. We establish the inequality with arbitrary
fixed boundary spins, as needed for the annuli below.

\subsection{Gauge invariance and second variations}

Local changes of spin variables underlie the Ward identities used
in~\cite{DriesslerLandauPerez,AizenmanSimonWard}; we begin with the
second-variation identity on a finite graph with prescribed spins.

Let \(G=(V,E)\) be a finite graph, with one orientation \(e=(x,y)\)
chosen for each edge, and let \(0\le b_e\le\beta\). Fix a set
\(P\subseteq V\) of vertices and arbitrary spins \(s_x\in S^2\) for
\(x\in P\). At the remaining vertices, integrate against uniform
probability measure on \(S^2\), with normalized density proportional to
\(\exp(\sum_{e=(x,y)}b_e s_x^Ts_y)\). Write \(\langle\cdot\rangle\)
for this expectation.

Let \(a,b,c\in\mathfrak{so}(3)\) denote cross product with the first,
second, and third coordinate vectors, respectively. Thus
\([a,b]=c\), and their operator norms are one. For an edge field
\(A=(A_e)_{e\in E}\) of skew matrices, put
\begin{equation}
\label{eq:rotation-partition}
\begin{split}
H(A)&=\sum_{e=(x,y)}b_e s_x^T\bigl(\exp(A_e)-I\bigr)s_y,\\
W(A)&=\bigl\langle\exp H(A)\bigr\rangle.
\end{split}
\end{equation}
Thus \(W\) is the ratio of the twisted partition function to the
untwisted one. All derivatives below are evaluated at \(A=0\). Set
\(X=DH\) and \(H_k=D^kH\) for \(k=2,3,4\). Compactness of the
finite spin space permits differentiation under the integral. We extend
all multilinear derivatives complex linearly in their arguments; changes
of spin variables themselves will always use real rotations.

For a scalar function on vertices, define \(df(e)=f_y-f_x\).
All scalar norms use counting measure on the indicated vertices or
edges. For a matrix indexed by edges,
\(\|T\|_{\mathrm{HS}}^2=\sum_{e,l\in E}|T_{el}|^2\).
For later use, the local derivative coefficients are
\begin{equation}
\label{eq:rotation-local-coefficients}
J_e(g_1,\ldots,g_k)
=b_e s_x^T\left(\frac{1}{k!}\sum_{\pi\in\mathfrak S_k}
g_{\pi(1)}\cdots g_{\pi(k)}\right)s_y,
\qquad |J_e(g_1,\ldots,g_k)|\le\beta,
\end{equation}
where each \(g_r\) is one of \(a,b,c\). Indeed, every matrix product
in this average has operator norm at most one. Consequently
\begin{equation}
\label{eq:rotation-local-expansion}
H_k(v_1g_1,\ldots,v_kg_k)
=\sum_e J_e(g_1,\ldots,g_k)\prod_{r=1}^k v_r(e).
\end{equation}
The coefficients depend on the spins and the generators, but not on the
scalar test functions. The same formula with \(k=1\) describes \(X\).

Define the complex bilinear form
\begin{equation}
\label{eq:rotation-R-definition}
R(h,k)=-D^2W(hc,kc)
\end{equation}
on scalar edge forms. It is symmetric, but is not assumed to be positive
semidefinite.

\begin{lemma}[Gradient estimates]\label{lem:gradient}
If \(f:V\to\mathbb C\) vanishes on \(P\), then, for
\(g\in\{a,b,c\}\),
\begin{equation}
\label{eq:rotation-gradient-bound}
\bigl\langle |X((df)g)|^2\bigr\rangle
\le\beta\|df\|_2^2.
\end{equation}
For every complex edge form \(h\),
\begin{equation}
\label{eq:rotation-R-bound}
R(h,\bar h)
=-\bigl\langle H_2(hc,\bar h c)+|X(hc)|^2\bigr\rangle
\le\beta\|h\|_2^2.
\end{equation}
Moreover, adding \(df\) in either argument of \(R\) does not change
its value whenever \(f\) vanishes on \(P\).
\end{lemma}

\begin{proof}
For real \(f\), make the change of variables
\(s_x^{\mathrm{old}}=\exp(-tf_xg)s_x^{\mathrm{new}}\) at every
integrated vertex. The fixed spins are unchanged because \(f|_P=0\).
Commutation of rotations about the same axis gives
\(W(t(df)g)=1\). Its second derivative at zero is
\[
0=\bigl\langle H_2((df)g,(df)g)+X((df)g)^2\bigr\rangle.
\]
The coefficient bound in \eqref{eq:rotation-local-coefficients} proves
\eqref{eq:rotation-gradient-bound} for real \(f\). For
\(f=f_1+if_2\), with \(f_1,f_2\) real, the squared modulus of
\(X((df)g)\) is the sum of the squares for \(f_1\) and \(f_2\).
Adding their bounds proves the complex case with the same constant.

Differentiating \eqref{eq:rotation-partition} twice gives the equality
in \eqref{eq:rotation-R-bound}. The local term is bounded in absolute
value by \(\beta\|h\|_2^2\), and the other term is nonnegative
before its minus sign, proving the inequality. Finally, the same change
of variables gives
\(W((h+df)c)=W(hc)\) for real \(h,f\). Differentiation in independent
scalar multiples of \(h\) and \(df\), followed by complex linear
extension, gives \(R(h,df)=R(df,h)=0\). This proves the last assertion.
\end{proof}

\subsection{The family inequality}

For vertex functions \(f,g\), write
\begin{equation}
\label{eq:rotation-S-definition}
S(f,g)_e=\tfrac12(f_xg_y-g_xf_y),\qquad e=(x,y).
\end{equation}
This form is bilinear and antisymmetric in \(f,g\).

\begin{proposition}[Fourth-order rotation inequality]\label{prop:fourth}
Let \((f_i)_{i\in I}\) be a finite family of complex vertex functions
vanishing on \(P\), and put
\[
u_i=df_i,\qquad
T_{el}=\sum_{i\in I}u_i(e)\overline{u_i(l)},\qquad
Q=\sum_{i\in I}S(f_i,\bar f_i).
\]
Then
\begin{equation}
\label{eq:rotation-family-bound}
\begin{split}
R(Q,\bar Q)
&\le K_{\mathrm{rot}}(\beta)\|T\|_{\mathrm{HS}}^2\\
&\quad+\beta\sum_{i,j\in I}
\min\bigl\{\|f_i\|_\infty^2\|df_j\|_2^2,
             \|f_j\|_\infty^2\|df_i\|_2^2\bigr\},
\end{split}
\end{equation}
where
\begin{equation}
\label{eq:rotation-explicit-constant}
K_{\mathrm{rot}}(\beta)=6\beta^2+4\beta^{3/2}+\beta.
\end{equation}
In particular, the constants are independent of the graph, the family
size, and the values of the fixed spins.
\end{proposition}

\begin{proof}
We first derive an exact identity relating the response \(R\) to a fourth
derivative of \(W\), and sum it over the family.  Gradient invariance will
control the resulting cross terms; a nonnegative square will supply the
remaining fourth-derivative lower bound.

\medskip\noindent
\emph{The noncommuting gauge identity.}
Let \(F_x\) be a real matrix potential in the span of \(a,b\), zero
on \(P\). The substitution
\(s_x^{\mathrm{old}}=\exp(-F_x)s_x^{\mathrm{new}}\) gives
\begin{equation}
\label{eq:rotation-nonabelian-gauge}
W(dF)=W(B),\qquad
B_e=\log\bigl(\exp(F_x)\exp(F_y-F_x)\exp(-F_y)\bigr)
\end{equation}
for \(F\) sufficiently small. The product inside the logarithm is
orthogonal and near the identity; hence its local logarithm is skew
symmetric. The logarithm also commutes with conjugation. Multiplying
the exponential series to degree two, and then taking the logarithm,
gives
\begin{equation}
\label{eq:rotation-B-quadratic}
B_e=\tfrac12[F_x,F_y]+O(F^3),\qquad DB_e(0)=0.
\end{equation}
For directions \(fa\) and \(gb\), the mixed second derivative of
\(B_e\) is therefore \(S(f,g)_e c\); it is zero for two directions
using the same generator.

We also need a fourth-order cancellation. Take four independent
parameters whose matrix directions have generator pattern \((a,b,a,b)\),
with arbitrary spatial potentials. The corresponding mixed fourth
derivative of \(B_e\) is zero. To see this, conjugate by the half-turn
about the first axis. This fixes \(a\) and negates \(b\), so the
coefficient involving all four parameters is fixed by conjugation:
its two \(b\)-directions contribute two minus signs. Conjugation by
the half-turn about the second axis likewise fixes this coefficient.
On \(\mathfrak{so}(3)\), the fixed spaces of these two conjugations
are respectively \(\mathbb R a\) and \(\mathbb R b\), with zero
intersection. Thus
\begin{equation}
\label{eq:rotation-B-fourth}
D_{1234}B_e=0
\quad\text{for the generator pattern }(a,b,a,b).
\end{equation}
Both \eqref{eq:rotation-B-quadratic} and
\eqref{eq:rotation-B-fourth} extend complex multilinearly. No complex
change of spin variables is involved in this extension.

For fixed \(i,j\), take the four edge directions
\begin{equation}
\label{eq:rotation-four-directions}
A_1=u_i a,\qquad A_2=\bar u_i b,\qquad
A_3=\bar u_j a,\qquad A_4=u_j b.
\end{equation}
The fourth derivative of the composition in
\eqref{eq:rotation-nonabelian-gauge} has no term containing \(DB\),
by \eqref{eq:rotation-B-quadratic}. Its term \(DW[D_{1234}B]\)
also vanishes, by \eqref{eq:rotation-B-fourth}. Only pairings of two
second derivatives remain. The pairing \(13|24\) is zero because
both pairs have the same generator. For the other two pairings, put
\(S_i=S(f_i,\bar f_i)\) and \(S_{ij}=S(f_i,f_j)\). Their arguments
are
\[
12|34:\quad(S_i c,\overline{S_j}c),\qquad
14|23:\quad(S_{ij}c,-\overline{S_{ij}}c).
\]
Here \(S(\bar f_j,f_j)=\overline{S_j}\) and
\(S(\bar f_j,\bar f_i)=-\overline{S_{ij}}\). Each labeled pairing
occurs once in the composition rule, so
\begin{equation}
\label{eq:rotation-fourth-identity}
D^4W(A_1,A_2,A_3,A_4)
=-R(S_i,\overline{S_j})+R(S_{ij},\overline{S_{ij}}).
\end{equation}
This also holds when some spatial functions coincide: the derivative
identity is multilinear in the four labeled directions.

Summing over \(i,j\) and rearranging gives the identity to be estimated:
\begin{equation}
\label{eq:rotation-response-balance}
R(Q,\bar Q)
=\sum_{i,j}R(S_{ij},\overline{S_{ij}})
 -\sum_{i,j}D^4W(u_i a,\bar u_i b,\bar u_j a,u_j b).
\end{equation}
We bound the first sum from above using gradient invariance, and the
second from below by retaining a nonnegative square.

\medskip\noindent
\emph{The cross terms.}
For functions \(f,g\) vanishing on \(P\), the product \(fg\)
also vanishes there, and the edge identities
\begin{equation}
\label{eq:rotation-gradient-adjustment}
\begin{split}
S(f,g)+\tfrac12d(fg)&=\tfrac12(f_x+f_y)\,dg,\\
S(f,g)-\tfrac12d(fg)&=-\tfrac12(g_x+g_y)\,df
\end{split}
\end{equation}
hold on \(e=(x,y)\). Applying gradient invariance in both arguments
of \(R\), followed by \eqref{eq:rotation-R-bound}, gives
\begin{equation}
\label{eq:rotation-cross-bound}
R(S(f,g),\overline{S(f,g)})
\le\beta\min\{\|f\|_\infty^2\|dg\|_2^2,
                 \|g\|_\infty^2\|df\|_2^2\}.
\end{equation}

\medskip\noindent
\emph{The fourth-derivative lower bound.}
We claim that
\begin{equation}
\label{eq:rotation-fourth-lower}
\sum_{i,j}D^4W(u_i a,\bar u_i b,\bar u_j a,u_j b)
\ge-K_{\mathrm{rot}}(\beta)\|T\|_{\mathrm{HS}}^2.
\end{equation}
To enumerate its terms, set \(X_r=X(A_r)\), and write
\(H_{r_1\cdots r_k}=H_k(A_{r_1},\ldots,A_{r_k})\). Direct
differentiation of \(\exp H\) gives all fifteen set partitions:
\begin{equation}
\label{eq:rotation-fifteen-terms}
\begin{aligned}
D^4W(A_1,A_2,A_3,A_4)=\bigl\langle
&X_1X_2X_3X_4\\
&+H_{12}X_3X_4+H_{13}X_2X_4+H_{14}X_2X_3\\
&+H_{23}X_1X_4+H_{24}X_1X_3+H_{34}X_1X_2\\
&+H_{12}H_{34}+H_{13}H_{24}+H_{14}H_{23}\\
&+H_{123}X_4+H_{124}X_3+H_{134}X_2+H_{234}X_1\\
&+H_{1234}\bigr\rangle.
\end{aligned}
\end{equation}
The four terms comprising
\((X_1X_2+H_{12})(X_3X_4+H_{34})\), after summation over \(i,j\),
have expectation
\begin{equation}
\label{eq:rotation-positive-square}
\left\langle\left|\sum_i
\bigl[X(u_i a)X(\bar u_i b)+H_2(u_i a,\bar u_i b)\bigr]
\right|^2\right\rangle\ge0.
\end{equation}
We bound the other eleven terms in absolute value.

For an edge \(e\), let \(t_e(l)=T_{el}\) and
\(\rho_e=(\sum_l|T_{el}|^2)^{1/2}\), and abbreviate
\(X_g(v)=X(vg)\). Crucially,
\begin{equation}
\label{eq:rotation-row-gradient}
t_e=d\left(\sum_i u_i(e)\overline{f_i}\right).
\end{equation}
The potential on the right, and its conjugate, vanish on \(P\).
Lemma~\ref{lem:gradient} therefore gives
\begin{equation}
\label{eq:rotation-row-bound}
\langle|X_g(t_e)|^2\rangle\le\beta\rho_e^2,
\qquad
\langle|X_g(\bar t_e)|^2\rangle\le\beta\rho_e^2.
\end{equation}

The following table records every remaining contraction. Expand each
\(H_k\) using \eqref{eq:rotation-local-expansion}, sum over \(i,j\),
and multiply the last two columns. One then sums over \(e\), except
in the two \(H_2H_2\) rows, where one sums over \(e,l\).
\begin{center}
\renewcommand{\arraystretch}{1.18}
\begin{tabular}{@{}lll@{}}
\toprule
Term & Local coefficient & Contracted expression\\
\midrule
\(H_{13}X_2X_4\) & \(J_e(a,a)\)
  & \(X_b(t_e)X_b(\bar t_e)\)\\
\(H_{14}X_2X_3\) & \(J_e(a,b)\)
  & \(X_b(t_e)X_a(t_e)\)\\
\(H_{23}X_1X_4\) & \(J_e(b,a)\)
  & \(X_a(\bar t_e)X_b(\bar t_e)\)\\
\(H_{24}X_1X_3\) & \(J_e(b,b)\)
  & \(X_a(\bar t_e)X_a(t_e)\)\\
\midrule
\(H_{13}H_{24}\) & \(J_e(a,a)J_l(b,b)\)
  & \(T_{el}\overline{T_{el}}\)\\
\(H_{14}H_{23}\) & \(J_e(a,b)J_l(b,a)\)
  & \(T_{el}^{\,2}\)\\
\midrule
\(H_{123}X_4\) & \(J_e(a,b,a)\)
  & \(T_{ee}X_b(\bar t_e)\)\\
\(H_{124}X_3\) & \(J_e(a,b,b)\)
  & \(T_{ee}X_a(t_e)\)\\
\(H_{134}X_2\) & \(J_e(a,a,b)\)
  & \(T_{ee}X_b(t_e)\)\\
\(H_{234}X_1\) & \(J_e(b,a,b)\)
  & \(T_{ee}X_a(\bar t_e)\)\\
\midrule
\(H_{1234}\) & \(J_e(a,b,a,b)\) & \(T_{ee}^{\,2}\)\\
\bottomrule
\end{tabular}
\end{center}
For example, the first row follows from
\[
\begin{split}
&\sum_{i,j}u_i(e)\bar u_j(e)
             X_b(\bar u_i)X_b(u_j)\\
&\qquad=
X_b\left(\sum_i u_i(e)\bar u_i\right)
X_b\left(\sum_j\bar u_j(e)u_j\right)
=X_b(t_e)X_b(\bar t_e).
\end{split}
\]
The other rows follow by the same finite summation; their explicit
expressions distinguish the conjugated and unconjugated contractions.

For each of the first four rows, the pointwise coefficient bound and
Cauchy--Schwarz, followed by \eqref{eq:rotation-row-bound}, give an
absolute expectation at edge \(e\) of at most
\(\beta^2\rho_e^2\). The coefficient may depend on the spins;
its pointwise absolute value is bounded before applying
Cauchy--Schwarz. These four rows contribute at most
\(4\beta^2\|T\|_{\mathrm{HS}}^2\).
The next two rows contribute at most
\(2\beta^2\|T\|_{\mathrm{HS}}^2\), since both contractions have
absolute value \(|T_{el}|^2\).

For each of the four \(H_3X\) rows,
\eqref{eq:rotation-row-bound} bounds the total absolute expectation by
\[
\beta^{3/2}\sum_e T_{ee}\rho_e
\le\beta^{3/2}
\left(\sum_e T_{ee}^2\right)^{1/2}
\left(\sum_e\rho_e^2\right)^{1/2}
\le\beta^{3/2}\|T\|_{\mathrm{HS}}^2.
\]
Here \(T_{ee}=\sum_i|u_i(e)|^2\ge0\), and the squared diagonal
entries form part of the Hilbert--Schmidt sum. The last row is bounded
by \(\beta\sum_e T_{ee}^2\le\beta\|T\|_{\mathrm{HS}}^2\).

Finally, the sum on the left of \eqref{eq:rotation-fourth-lower} is
real: complex conjugation interchanges \(i,j\) and permutes the
arguments of the symmetric fourth derivative. Keeping the nonnegative
square \eqref{eq:rotation-positive-square}, and subtracting the bounds
for the remaining terms, proves \eqref{eq:rotation-fourth-lower} with
exactly \(K_{\mathrm{rot}}(\beta)\) from
\eqref{eq:rotation-explicit-constant}.

Combining \eqref{eq:rotation-cross-bound} and
\eqref{eq:rotation-fourth-lower} in
\eqref{eq:rotation-response-balance} proves
\eqref{eq:rotation-family-bound}. Only the upper bound and gradient
invariance of \(R\) were used; no positivity property of \(R\) is
required.
\end{proof}
